\chapter{The Laser Beam and the Echo}

\section{The beam spreads out}
\begin{kid}
Shine a flashlight at your hand: you see a small circle. Now shine it at a wall across the room: the circle is much bigger. A laser beam spreads too --- just very, very slowly. Far away, the laser ``dot'' can be as big as a coin, so it can touch the edge of a pole and the wall behind it at the same time.
\end{kid}

\fig[0.82]{footprint}{The footprint of a beam with 5 mm exit diameter and 0.3 mrad divergence.}

\begin{keyeq}{Beam footprint}
\[ D(\rho)\approx d_0+\rho\,\gamma \]
\begin{symbols}
$D$ & diameter of the laser spot on the target & m\\
$d_0$ & diameter of the beam as it leaves the sensor & m\\
$\rho$ & range & m\\
$\gamma$ & full-angle beam divergence (small-angle rule again!) & rad (often given in mrad = $10^{-3}$ rad)\\
\end{symbols}
\end{keyeq}

\begin{worked}{spot size at 50 m and 100 m}
\textbf{Given:} $d_0=5$ mm $=0.005$ m, $\gamma=0.3$ mrad $=0.0003$ rad.
\step{At 50 m} $D=0.005+50\times0.0003=0.005+0.015=0.020$ m $=2.0$ cm
\step{At 100 m} $D=0.005+100\times0.0003=0.035$ m $=3.5$ cm
\result{a 2--3.5 cm ``blur'' at mapping ranges. You cannot locate an edge more sharply than about half the spot size.}
\end{worked}

\section{How strong is the echo? The LiDAR range equation}
\begin{kid}
The farther away something is, the dimmer its echo. It gets dim fast: twice as far means only one quarter as bright. Shiny white things (a road sign) send back lots of light; black asphalt or wet surfaces send back very little. And if the light hits a surface at a slant, most of it bounces away like a ball skipping off water.
\end{kid}

\begin{keyeq}{Received power for an extended diffuse target}
\[ P_r = P_t\,\frac{\rho_{\text{refl}}\,\cos\theta_i}{\pi\,R^2}\,A_r\,\eta_{\text{sys}}\,\eta_{\text{atm}}^2 \]
\begin{symbols}
$P_r$ & power arriving at the detector & W\\
$P_t$ & power transmitted in the pulse & W\\
$\rho_{\text{refl}}$ & target reflectance (0 = black, 1 = perfect white) & ---\\
$\theta_i$ & incidence angle between beam and surface normal & rad\\
$R$ & range & m\\
$A_r$ & area of the receiving lens & m$^2$\\
$\eta_{\text{sys}}$ & optical efficiency of the sensor & ---\\
$\eta_{\text{atm}}$ & one-way atmospheric transmission (squared: out and back) & ---\\
$\pi$ & comes from a diffuse (Lambertian) surface scattering into a half-sphere & ---\\
\end{symbols}
\end{keyeq}

\fig[0.95]{power}{Left: the $1/R^2$ law. Right: the $\cos\theta_i$ incidence-angle law.}

\begin{worked}{comparing two targets}
\textbf{Given:} Target A: white concrete ($\rho_{\text{refl}}=0.6$) at 10 m, head-on. Target B: asphalt ($\rho_{\text{refl}}=0.1$) at 50 m, hit at 60°. Everything else equal.
\step{Step 1 --- ratio} $\dfrac{P_B}{P_A}=\dfrac{0.1}{0.6}\times\dfrac{\cos60^\circ}{\cos0^\circ}\times\dfrac{10^2}{50^2}$
\step{Step 2} $=0.1667\times0.5\times0.04=0.00333$
\result{B returns about 1/300 of A's power. Weak echoes have noisier ranges, which is why range precision gets worse with distance and grazing angle --- and why road surfaces far from the vehicle are the noisiest part of a mobile scan.}
\end{worked}

\section{Intensity}
Every point carries an \textbf{intensity} value, a number related to $P_r$. Because of the $R^2$ and $\cos\theta_i$ effects, raw intensity of the same paint changes with range and angle; software can \emph{normalize} it. Intensity makes road markings, signs and retro-reflective targets jump out of the cloud, which is handy for extracting line work and for finding control targets.

\section{Wavelength and eye safety}
\begin{eng}
Most automotive and mobile scanners use \textbf{905 nm} (near-infrared, cheap silicon detectors) or \textbf{1550 nm} (absorbed by the eye's cornea and lens before reaching the retina, so far more power may be used eye-safely --- giving longer range). All survey scanners used in public must be \textbf{Class 1} eye-safe under IEC 60825-1. Water strongly absorbs 1550 nm, so wet surfaces and rain reduce returns more at that wavelength; green (532 nm) lasers are used for bathymetric LiDAR because green light penetrates water.
\end{eng}

\section{What can LiDAR see through?}
\begin{kid}
A laser is light, so it acts like a flashlight beam, not like X-ray vision. It cannot see \emph{through} a wall, a tree trunk, or the ground. It \emph{can} slip through holes --- the gaps between leaves, the openings in a fence --- and reach whatever is behind them. It passes through clear glass the way sunlight does. And it is stopped almost right away by water.
\end{kid}

\fig[1.0]{penetration}{The five situations every LiDAR user meets. Orange dots are recorded returns.}

\begin{center}\small
\renewcommand{\arraystretch}{1.2}
\begin{tabularx}{\linewidth}{>{\raggedright\arraybackslash}p{3.0cm}>{\raggedright\arraybackslash}p{2.2cm}X}\toprule
\textbf{Material or condition} & \textbf{Does the laser get through?} & \textbf{What you get and what to do}\\\midrule
Tree canopy, shrubs & \textbf{Through the gaps only} & First returns on leaves, last returns on the ground where there are openings. Dense evergreen or summer canopy may leave no ground hits: collect in leaf-off season, slow down, add passes.\\
Tall grass, crops, mangroves & \textbf{Rarely} & Points sit on top of the vegetation, not the soil. Heights are biased high; take ground shots by conventional methods.\\
Chain-link fence, mesh, railings & \textbf{Partly} (gaps) & Some points on the fence, some behind it; the fence appears as a sparse ghost surface.\\
Clear glass windows & \textbf{Mostly through} & Few returns from the glass itself (except head-on); you see the room behind. Reflections can create ``mirror-world'' ghost points.\\
Tinted, frosted or coated glass & \textbf{Partly / varies} & Unpredictable; check the cloud.\\
Mirrors, polished metal, chrome & \textbf{No --- reflects away} & Dropouts or ghost points placed at false ranges behind the surface.\\
Water (865--1550 nm) & \textbf{No} & Absorbed within centimeters; calm water often returns nothing except straight down. Only green (532 nm) bathymetric LiDAR maps the bottom, in clear water.\\
Wet pavement, puddles & \textbf{No} & Weak or missing returns; specular reflection; schedule collection after surfaces dry.\\
Rain, fog, dust, smoke, steam & \textbf{Partly} & Range shortens; ``floating'' false points near the sensor; heavy fog can block completely. Filter them out (Module 18).\\
Falling snow & \textbf{Partly} & Noise points in the air; snow cover hides the ground surface.\\
Walls, soil, pavement, people, vehicles & \textbf{No} & Opaque: anything behind is in shadow. Plan extra passes or positions to see behind objects.\\
Underground utilities & \textbf{No} & LiDAR maps only visible surfaces (manhole lids, valve covers). Use ground-penetrating radar or locating for what is buried.\\\bottomrule
\end{tabularx}\end{center}

\begin{bigidea}
LiDAR never ``penetrates'' solid material. It reaches hidden ground only through openings, which is why multiple returns (Module 3) and many viewing angles matter more than laser power.
\end{bigidea}

\section{Reflectivity of common surfaces}
\fig[0.86]{reflectivity}{Typical reflectivity of common surfaces at the near-infrared wavelengths used by most mobile scanners (865--905 nm). Values vary with moisture, age, roughness, color and incidence angle.}

\begin{center}\small
\renewcommand{\arraystretch}{1.18}
\begin{tabularx}{\linewidth}{>{\raggedright\arraybackslash}p{4.2cm}>{\centering\arraybackslash}p{2.3cm}X}\toprule
\textbf{Surface} & \textbf{Typical reflectivity} & \textbf{What it means in the field}\\\midrule
Clear, calm water & 1--5\% & often no return at all except straight below the sensor\\
Black tire, rubber & 2--5\% & among the darkest targets; shortest detection range\\
Black or dark car paint & 3--8\% & cars may appear with holes; glossy paint can also mirror\\
Fresh asphalt & 4--8\% & dark: road returns thin out first at long range\\
Wet soil, wet asphalt & 4--12\% & water darkens surfaces in the near infrared\\
Dark rock (basalt) & 5--15\% & dark\\
Aged, gray asphalt & 10--20\% & brighter as it weathers\\
Dry soil & 15--35\% & medium\\
Tree bark, wood & 20--45\% & medium\\
Concrete & 25--50\% & good, consistent target\\
Dry sand & 30--55\% & bright\\
Grass, green leaves & 40--60\% & bright in the near infrared (healthy plants reflect strongly there)\\
White paint, road striping & 60--85\% & very bright: striping stands out in intensity images\\
Fresh snow & 75--95\% at 865--905 nm & very bright at 905 nm --- but much darker at 1550 nm\\
Retroreflective signs and tape & ``over 100\%'' & sends light straight back; can saturate the receiver and bias the range (Module 15)\\\bottomrule
\end{tabularx}\end{center}
{\footnotesize Values are typical ranges compiled for teaching; they are not specifications. Sensors that report \emph{calibrated reflectivity} (such as the Ouster OS0) let you measure your own surfaces.}

\begin{worked}{how far can I see dark asphalt?}
\textbf{Given:} the OS0 detects a 10\% target to 35 m (Module 15). Received power falls as $1/R^2$ (Section 4.2), so for the same detection threshold the maximum range grows as the \emph{square root} of reflectivity.
\step{Fresh asphalt, 5\%} $35\times\sqrt{5/10}=35\times0.71=25$ m
\step{Black car, 4\%} $35\times\sqrt{0.4}=22$ m
\step{Concrete, 40\%} $35\times\sqrt{4}=70$ m (the datasheet's 75 m at 80\% shows the rule flattening out in bright sunlight)
\result{plan the usable swath on the \emph{darkest} surface you must map: on new asphalt, expect road coverage to end near 25 m, long before the brochure range.}
\end{worked}

\section{Range precision in practice}
\begin{center}\small
\renewcommand{\arraystretch}{1.2}
\begin{tabularx}{\linewidth}{Xl}\toprule
\textbf{Situation} & \textbf{Typical range noise (1$\sigma$)}\\\midrule
Good mobile unit, bright hard target, under 30 m (100 ft) & 1--3 cm (0.03--0.10 ft)\\
Same unit, dark or grazing target, 50--100 m (165--330 ft) & 5--15 cm (0.16--0.49 ft)\\
Survey-grade mobile mapping head & 3--8 mm (0.01--0.03 ft)\\
Terrestrial phase scanner, short range & 1--3 mm (0.003--0.010 ft)\\\bottomrule
\end{tabularx}\end{center}
These are order-of-magnitude illustrations; always read the specification sheet and, more importantly, \emph{test on your own control}.

\begin{tryit}
\textbf{4.1} A beam has $d_0=8$ mm and $\gamma=0.5$ mrad. How wide is the spot at 40 m?\\
\textbf{4.2} By what factor does echo power drop when a target moves from 20 m to 60 m?\\
\textbf{4.3} (Simple) Which gives a brighter echo: a white sign or a black tire? Why?\\
\textbf{4.4} (Simple) Which of these can a mobile LiDAR measure the ground behind: a leafy tree, a brick wall, a glass storefront, a pond?
\end{tryit}

\sources{Baltsavias (1999), \emph{ISPRS J.} 54 (range equation, footprint); IEC 60825-1:2014 (laser classes); Ouster OS0 Rev8 datasheet (calibrated reflectivity, detection ranges at 10\% and 80\%); reflectivity values are typical teaching ranges --- measured spectra are available in the USGS Spectral Library Version 7 (Kokaly \emph{et al.}, 2017).}
